% =====================================================================
%  Informe del Proyecto 3 - Grupo 3 - Desarrollo de Videojuegos (ESPE)
%  Tres en Raya con IA: backend Go + frontend Angular
%  Compilar:  pdflatex Grupo3-InformeProy.tex  (dos veces, para el índice)
% =====================================================================
\documentclass[12pt,a4paper]{article}

\usepackage[utf8]{inputenc}
\usepackage[T1]{fontenc}
\usepackage[spanish,es-tabla,es-noshorthands]{babel}
\usepackage{lmodern}
\usepackage[margin=2.5cm]{geometry}
\usepackage{graphicx}
\usepackage{xcolor}
\usepackage{booktabs}
\usepackage{longtable}
\usepackage{array}
\usepackage{enumitem}
\usepackage{listings}
\usepackage{float}
\usepackage{caption}
\usepackage{fancyhdr}
\usepackage{tikz}
\usetikzlibrary{shapes.geometric,shapes.multipart,arrows.meta,positioning,calc,babel}
\usepackage[hidelinks]{hyperref}

\graphicspath{{Img/}}

% ---------- colores y utilidades ----------
\definecolor{espe}{RGB}{0,102,51}
\definecolor{xcol}{RGB}{8,145,178}
\definecolor{ocol}{RGB}{219,39,119}
\definecolor{pend}{RGB}{180,120,0}
\definecolor{codebg}{RGB}{245,245,250}

% Marca para los apartados que completan los compañeros
\newcommand{\pendiente}[1]{\par\noindent\fcolorbox{pend}{pend!8}{\parbox{\dimexpr\linewidth-2\fboxsep-2\fboxrule}{\small\textbf{\textcolor{pend}{Pendiente:}} #1}}\par\medskip}

\lstdefinelanguage{Go}{
  morekeywords={package,import,func,return,if,else,for,range,switch,case,default,type,struct,var,const,interface,map,go,defer,true,false,nil},
  sensitive=true, morecomment=[l]{//}, morestring=[b]"}
\lstset{basicstyle=\ttfamily\footnotesize, backgroundcolor=\color{codebg},
  keywordstyle=\color{xcol}\bfseries, commentstyle=\color{gray}\itshape,
  stringstyle=\color{ocol}, frame=single, rulecolor=\color{gray!40},
  breaklines=true, columns=fullflexible, showstringspaces=false,
  literate={á}{{\'a}}1 {é}{{\'e}}1 {í}{{\'i}}1 {ó}{{\'o}}1 {ú}{{\'u}}1 {ñ}{{\~n}}1}

\pagestyle{fancy}
\fancyhf{}
\fancyhead[L]{\small Desarrollo de Videojuegos -- Unidad 1}
\fancyhead[R]{\small Grupo N.º 3}
\fancyfoot[C]{\thepage}
\renewcommand{\headrulewidth}{0.4pt}
\setlength{\headheight}{15pt}

\begin{document}
\sloppy

% =====================================================================
%  CARÁTULA
% =====================================================================
\begin{titlepage}
  \centering
  \includegraphics[width=0.55\textwidth]{logoINFESPE.png}\par\vspace{1cm}
  {\Large\bfseries UNIVERSIDAD DE LAS FUERZAS ARMADAS -- ESPE\par}
  \vspace{0.3cm}
  {\large DEPARTAMENTO DE CIENCIAS DE LA COMPUTACIÓN\par}
  \vspace{0.2cm}
  {\large CARRERA DE INGENIERÍA EN SOFTWARE\par}
  \vspace{1.2cm}
  {\large\bfseries ASIGNATURA: DESARROLLO DE VIDEOJUEGOS\par}
  \vspace{0.3cm}
  {\large UNIDAD 1\par}
  \vspace{1.2cm}
  {\color{espe}\rule{\linewidth}{1pt}}\par\vspace{0.3cm}
  {\LARGE\bfseries Tres en Raya con Inteligencia Artificial\par}
  \vspace{0.2cm}
  {\large Proyecto 3: Tic-Tac-Toe (UnbeatableAI + MediumAI) -- Go (Web) + Angular\par}
  \vspace{0.3cm}
  {\color{espe}\rule{\linewidth}{1pt}}\par
  \vspace{1.2cm}
  {\large DOCENTE: ING. CÉSAR VILLACÍS SILVA, MSC.\par}
  \vspace{0.8cm}
  {\large\bfseries ESTUDIANTES:\par}
  \vspace{0.2cm}
  {\large Diana Guerra\par Paul Moreno\par Germán Cáceres\par}
  \vspace{0.8cm}
  {\large\bfseries GRUPO N.º 3\par}
  \vfill
  {\large FECHA: NOVIEMBRE 2026\par}
\end{titlepage}

\tableofcontents
\newpage

% =====================================================================
\section{Proyectos Propuestos del Capítulo 1}
% =====================================================================
El docente propuso implementar el juego del Tres en Raya con niveles de
dificultad en distintas tecnologías. Al Grupo 3 le corresponden los
proyectos web en Go y JavaScript:

\begin{table}[H]
\centering
\caption{Proyecto 3 -- Grupo 3}
\begin{tabular}{clll}
\toprule
\textbf{N.º} & \textbf{Variante de IA} & \textbf{Tecnología} & \textbf{Estado} \\
\midrule
1 & UnbeatableAI & Go (Web) + Angular & \textbf{Implementado} \\
2 & MediumAI & Go (Web) + Angular & \textbf{Implementado} \\
3 & UnbeatableAI & JavaScript (Web) & -- \\
4 & MediumAI & JavaScript (Web) & -- \\
5 & UnbeatableAI $\rightarrow$ MinMax & Go (Web) + Angular & Apartado preparado \\
6 & UnbeatableAI $\rightarrow$ MinMax & JavaScript (Web) & -- \\
\bottomrule
\end{tabular}
\end{table}

El equipo eligió \textbf{Go} para el backend y \textbf{Angular} para el
frontend. Por especificación, la lógica del juego es \emph{exactamente la
misma} que la de la versión oficial del docente
(\texttt{Tic-Tac-Toe-Game-Windows-Professor}); solo cambia la tecnología.
Esa versión trae una aplicación de Windows Forms por variante:

\begin{table}[H]
\centering
\caption{Versiones de referencia del docente}
\small
\begin{tabular}{lll}
\toprule
\textbf{Variante} & \textbf{Carpeta} & \textbf{Clase}\\
\midrule
UnbeatableAI & \texttt{Tic-Tac-Toe-Versión-IA-Fuerte} & \texttt{UnbeatableAIController}\\
MediumAI & \texttt{Tic-Tac-Toe-Versión-IA-Debil} & \texttt{MediumAIController}\\
MinMax & \texttt{Tic-Tac-Toe-Version-MiniMax} & \texttt{UnbeatableAIController}\textsuperscript{*}\\
\bottomrule
\end{tabular}\\[2pt]
{\footnotesize \textsuperscript{*}Misma clase, pero implementada con el algoritmo Minimax en lugar del método numérico.}
\end{table}

% =====================================================================
\section{Especificación de Requerimientos (IEEE 830)}
% =====================================================================

\subsection{Introducción}

\subsubsection{Propósito}
Este documento especifica los requerimientos del videojuego educativo
\emph{Tres en Raya con IA} en su versión web. Va dirigido al docente de la
asignatura y a los integrantes del grupo.

\subsubsection{Alcance}
El sistema permite que un usuario juegue al Tres en Raya contra un
\emph{Jugador Controlado por la Computadora} (JCC) desde un navegador. La
JCC decide sus jugadas con el método numérico heurístico de la referencia
(pesos usuario = 1, JCC = 3, vacío = 0). Los niveles implementados son \textbf{UnbeatableAI} y \textbf{MediumAI};
\textbf{UnbeatableAI $\rightarrow$ MinMax} permanece preparado.

\subsubsection{Definiciones, acrónimos y abreviaturas}
\begin{description}[leftmargin=2.6cm,style=nextline]
  \item[JCC] Jugador Controlado por la Computadora (marca ``O'').
  \item[Usuario] Jugador humano (marca ``X''); siempre empieza.
  \item[mArrayMoves] Vector $V_k$ de 9 posiciones con los pesos 0/1/3.
  \item[API REST] Interfaz HTTP/JSON entre el frontend y el backend.
  \item[SPA] \emph{Single Page Application}, aplicación web de una sola página.
\end{description}

\subsubsection{Referencias}
\begin{itemize}
  \item IEEE Std 830-1998, \emph{Recommended Practice for Software Requirements Specifications}.
  \item Villacís, C. et al., \emph{Optimización del Juego del Tres en Raya con
        Niveles de Dificultad, Utilizando Técnicas Heurísticas de IA}, ESPE.
  \item Material de clase: \emph{Caso 1}, \emph{Caso 2} y \emph{Caso 3 TTT}.
  \item Código del docente: \texttt{Tic-Tac-Toe-Game-Windows-Professor} (IA fuerte, IA débil y MiniMax).
\end{itemize}

\subsection{Descripción general}

\subsubsection{Perspectiva del producto}
Es una aplicación cliente--servidor. El \textbf{backend en Go} contiene
\emph{toda} la lógica del juego (port 1:1 de las clases C\#) y la expone como
API REST. El \textbf{frontend en Angular} solo dibuja el tablero, envía los
clics y anima los eventos que devuelve el servidor.

\subsubsection{Funciones del producto}
Seleccionar modo de juego, iniciar una partida nueva, marcar una casilla,
recibir la jugada de la JCC, conocer el resultado (``¡Ganó X!'', ``¡Ganó O!'' o
``¡Empate!''), ver el marcador y ver cómo razona la IA (pesos y sumas por línea).

\subsubsection{Características de los usuarios}
Niños de 7 a 11 años (público objetivo del juego original) y estudiantes.
No se requiere conocimiento técnico: basta con hacer clic o usar las teclas
1--9.

\subsubsection{Restricciones}
\begin{itemize}
  \item La lógica del juego no puede modificarse respecto al original.
  \item Backend en Go~1.23 o superior, solo con la librería estándar.
  \item Frontend en Angular~22 con navegadores modernos (Chrome, Edge, Firefox).
\end{itemize}

\subsubsection{Suposiciones y dependencias}
El servidor Go está disponible en el puerto 8080; en desarrollo, Angular
redirige \texttt{/api} mediante \texttt{proxy.conf.json}.

\subsection{Requerimientos específicos}

\subsubsection{Requerimientos funcionales}
\begin{longtable}{p{1.3cm}p{9.2cm}p{2.6cm}}
\caption{Requerimientos funcionales}\\
\toprule
\textbf{ID} & \textbf{Descripción} & \textbf{Prioridad}\\
\midrule\endfirsthead
\toprule
\textbf{ID} & \textbf{Descripción} & \textbf{Prioridad}\\
\midrule\endhead
RF01 & El sistema mostrará un tablero de 3$\times$3 casillas numeradas 0--8 (TabIndex). & Alta\\
RF02 & El usuario podrá seleccionar el modo de juego; los modos no implementados se mostrarán como ``en construcción''. & Alta\\
RF03 & El usuario siempre inicia con la marca ``X''. & Alta\\
RF04 & Al marcar una casilla libre, el sistema registrará peso 1 en \texttt{mArrayMoves}; una casilla ocupada se ignora. & Alta\\
RF05 & En su primera jugada la JCC tomará el centro (4) si está libre; si no, una esquina aleatoria (0, 2, 6, 8). En las jugadas siguientes, si el centro está libre, lo toma. & Alta\\
RF06 & La JCC aplicará el método numérico: suma 6 $\Rightarrow$ gana; suma 2 $\Rightarrow$ bloquea; suma 5 en diagonal $\Rightarrow$ caso trivial; dos lados adyacentes del usuario $\Rightarrow$ esquina intermedia (orden 1+3$\to$0, 1+5$\to$2, 3+7$\to$6, 5+7$\to$8); en otro caso, casilla aleatoria. & Alta\\
RF07 & Si la JCC está en el centro y la regla de suma 5 aplica, jugará un lado aleatorio entre 1, 5 y 7 (comportamiento del código del docente, ver sección~\ref{sec:hallazgos}). & Alta\\
RF08 & El sistema verificará ganador recorriendo las 8 líneas (filas, columnas, diagonales); con 9 jugadas sin ganador declara empate. Mensajes: ``¡Ganó X!'', ``¡Ganó O!'', ``¡Empate!''. & Alta\\
RF09 & Al terminar la partida el tablero se deshabilitará (equivalente a \texttt{grbTicTacToe.Enabled = false}). & Alta\\
RF10 & El usuario podrá iniciar un juego nuevo en cualquier momento. & Media\\
RF11 & El sistema mostrará un marcador de partidas ganadas, empatadas y perdidas por modo. & Media\\
RF12 & El sistema mostrará los pesos de \texttt{mArrayMoves} y la suma de cada línea (inspector de IA). & Baja\\
\bottomrule
\end{longtable}

\subsubsection{Requerimientos no funcionales}
\begin{longtable}{p{1.3cm}p{2.6cm}p{9.2cm}}
\caption{Requerimientos no funcionales}\\
\toprule
\textbf{ID} & \textbf{Tipo} & \textbf{Descripción}\\
\midrule\endfirsthead
\toprule
\textbf{ID} & \textbf{Tipo} & \textbf{Descripción}\\
\midrule\endhead
RNF01 & Fidelidad & La lógica de la IA será idéntica a \texttt{UnbeatableAIController.cs} (versión IA Fuerte del docente).\\
RNF02 & Rendimiento & El servidor responderá a cada jugada en menos de 50 ms.\\
RNF03 & Usabilidad & Interfaz animada, con sonido opcional y soporte de teclado (1--9 y teclado numérico).\\
RNF04 & Accesibilidad & Respeta \texttt{prefers-reduced-motion} y usa etiquetas ARIA en el tablero.\\
RNF05 & Portabilidad & Funciona en navegadores modernos y es adaptable a celular, tablet y escritorio.\\
RNF06 & Mantenibilidad & Cada modo vive en su propio archivo y cumple una interfaz común (\texttt{Controller}).\\
RNF07 & Fiabilidad & Ninguna petición puede bloquear el servidor (protección ante bucles infinitos del original).\\
RNF08 & Concurrencia & Varias partidas simultáneas, aisladas por sesión, con exclusión mutua.\\
RNF09 & Calidad & Pruebas automatizadas con \texttt{go test} para los Casos 1, 2 y 3.\\
\bottomrule
\end{longtable}

% =====================================================================
\section{Modelo de Inteligencia Artificial}
% =====================================================================
El tablero se representa como el vector $V_k$, con $k = 0,\dots,8$, en el
arreglo \texttt{mArrayMoves}. Cada casilla vale 0 (vacía), 1 (usuario) o
3 (JCC). Para cada una de las 8 líneas se calcula la suma de sus tres
casillas:

\begin{table}[H]
\centering
\caption{Interpretación de la suma por línea}
\begin{tabular}{cll}
\toprule
\textbf{Suma} & \textbf{Composición} & \textbf{Acción de la JCC}\\
\midrule
6 & 3 + 3 + 0 & Completar la línea y \textbf{ganar} (Caso 2)\\
2 & 1 + 1 + 0 & \textbf{Bloquear} al usuario (Caso 1)\\
5 (diagonal) & 1 + 3 + 1 & Caso trivial: esquina libre o lado aleatorio (Caso 3)\\
\bottomrule
\end{tabular}
\end{table}

La regla de bloqueo en Go conserva el mismo orden de evaluación que el C\#:

\begin{lstlisting}[language=Go]
// pick devuelve el primer casillero vacío de la línea (a, b, d).
func (c *UnbeatableAIController) pick(a, b, d int) int {
	index := 0
	if c.ArrayMoves[a] == 0 { index = a
	} else if c.ArrayMoves[b] == 0 { index = b
	} else if c.ArrayMoves[d] == 0 { index = d }
	c.heuristicSum = c.ArrayMoves[a] + c.ArrayMoves[b] + c.ArrayMoves[d]
	return index
}
\end{lstlisting}

% =====================================================================

\subsection{Variante MediumAI}
\texttt{MediumAIController} porta la versi\'on de IA d\'ebil del docente y
conserva sus pesos: 0 para una casilla vac\'ia, 1 para X y 3 para O. El turno
humano siempre inicia con X. Despu\'es de cada jugada v\'alida, el controlador
comprueba las ocho l\'ineas ganadoras y el empate; al terminar, registra el
mensaje del resultado y deshabilita el tablero.

La JCC (O) aplica las reglas en este orden:
\begin{enumerate}
  \item Si el centro (\texttt{CenterIndex = 4}) est\'a vac\'io, lo elige de inmediato.
  \item Si el centro est\'a ocupado, recorre filas, columnas y diagonales en el
        orden de \texttt{WinningLines}. Primero busca una suma de 6
        ($3+3+0$) para completar una l\'inea y ganar.
  \item Si no hay una suma de 6, busca una suma de 2 ($1+1+0$) para bloquear
        la primera amenaza de X, eligiendo la primera casilla vac\'ia de esa l\'inea.
  \item Si ninguna regla aplica, elige al azar una casilla vac\'ia.
\end{enumerate}

La prioridad global de suma 6 antes de suma 2 permite que O gane cuando tiene
una jugada ganadora disponible, aunque exista una amenaza de X. Esta variante
no incluye las reglas de suma 5 ni la estrategia de esquinas y lados de
\texttt{UnbeatableAI}. El backend expone los pesos por medio de
\texttt{Weights()} y comunica jugadas y resultados como eventos del tablero;
la interfaz Angular los presenta sin decidir movimientos.

\section{Diagramas}
% =====================================================================

\subsection{Diagrama de casos de uso}
\begin{figure}[H]
\centering
\begin{tikzpicture}[
  uc/.style={ellipse,draw,fill=xcol!8,minimum width=4.4cm,minimum height=0.9cm,font=\small,align=center},
  actor/.style={font=\small,align=center},
  >=Stealth]
  \node[draw,rounded corners,minimum width=6.2cm,minimum height=10.4cm] (sys) at (0,0) {};
  \node[font=\bfseries\small,anchor=south] at (sys.north) {Tres en Raya IA (Web)};
  \node[uc] (u1) at (0,4.3)  {Seleccionar modo de juego};
  \node[uc] (u2) at (0,3.0)  {Iniciar juego nuevo};
  \node[uc] (u3) at (0,1.7)  {Marcar casilla};
  \node[uc] (u4) at (0,0.0)  {Calcular jugada de la JCC};
  \node[uc] (u5) at (0,-1.7) {Verificar ganador};
  \node[uc] (u6) at (0,-3.1) {Consultar marcador};
  \node[uc] (u7) at (0,-4.4) {Ver inspector de IA};
  \node[actor] (usr) at (-5.5,0.5) {\Huge\textcolor{xcol}{$\bigcirc$}\\[-4pt]\rule{1pt}{1cm}\\Usuario};
  \node[actor] (jcc) at (5.5,-0.8) {\Huge\textcolor{ocol}{$\bigcirc$}\\[-4pt]\rule{1pt}{1cm}\\JCC\\(Backend Go)};
  \foreach \u in {u1,u2,u3,u6,u7} \draw (usr.east) -- (\u.west);
  \draw (jcc.west) -- (u4.east);
  \draw (jcc.west) -- (u5.east);
  \draw[dashed,->] (u3) -- node[right,font=\scriptsize]{$\ll$include$\gg$} (u4);
  \draw[dashed,->] (u4) -- node[right,font=\scriptsize]{$\ll$include$\gg$} (u5);
\end{tikzpicture}
\caption{Diagrama de casos de uso}
\end{figure}

\subsection{Diagrama de secuencia: marcar casilla}
\begin{figure}[H]
\centering
\resizebox{\textwidth}{!}{%
\begin{tikzpicture}[font=\small,>=Stealth,
  obj/.style={draw,fill=xcol!10,minimum width=2.4cm,minimum height=0.9cm,align=center}]
  \node[obj] (a) at (0,0)    {:Usuario};
  \node[obj] (b) at (3.3,0)  {:Angular\\Board};
  \node[obj] (c) at (6.8,0)  {:GameStore};
  \node[obj] (d) at (10.3,0) {:API Go};
  \node[obj] (e) at (13.8,0) {:Form /\\UnbeatableAI};
  \foreach \n in {a,b,c,d,e} \draw[dashed,gray] (\n.south) -- ++(0,-10.2);
  \draw[->] (0,-1.2) -- node[above]{clic casilla $i$} (3.3,-1.2);
  \draw[->] (3.3,-1.9) -- node[above]{play($i$)} (6.8,-1.9);
  \draw[->] (6.8,-2.6) -- node[above]{POST /moves \{index\}} (10.3,-2.6);
  \draw[->] (10.3,-3.3) -- node[above]{Click($i$)} (13.8,-3.3);
  \draw[->] (13.8,-3.8) -- ++(0.8,0) -- node[right]{PlayUser + CheckWinner} ++(0,-0.6) -- (13.8,-4.4);
  \draw[->] (13.8,-5.1) -- ++(0.8,0) -- node[right]{PlayNPC (heurística)} ++(0,-0.6) -- (13.8,-5.7);
  \draw[->] (13.8,-6.4) -- ++(0.8,0) -- node[right]{CheckWinner} ++(0,-0.6) -- (13.8,-7.0);
  \draw[->,dashed] (13.8,-7.9) -- node[above]{eventos} (10.3,-7.9);
  \draw[->,dashed] (10.3,-8.5) -- node[above]{GameState JSON} (6.8,-8.5);
  \draw[->] (6.8,-9.1) -- node[above]{animar X, ``pensando'', O} (3.3,-9.1);
  \draw[->,dashed] (3.3,-9.8) -- node[above]{tablero + resultado} (0,-9.8);
\end{tikzpicture}}
\caption{Diagrama de secuencia del caso de uso ``Marcar casilla''. Si hay
ganador o empate, \texttt{Board.Enabled} pasa a \texttt{false} y el frontend
muestra el resultado con animación.}
\end{figure}

\subsection{Diagrama de clases (backend Go)}
\begin{figure}[H]
\centering
\resizebox{\textwidth}{!}{%
\begin{tikzpicture}[font=\scriptsize,>=Stealth,
  cls/.style={draw,rectangle split,rectangle split parts=3,fill=white,
              text width=4.8cm,align=left,inner sep=3pt},
  pend/.style={cls,draw=pend}]
  \node[draw,fill=gray!10,text width=3.4cm,align=center] (api) at (-5.4,0) {\textbf{api.Server}\\sesiones (map + mutex)\\rutas REST, CORS};
  \node[cls] (form) at (0,0) {\centering\textbf{Form} (frmTicTacToe)
    \nodepart{second}+GameLevel: int\\+ControlPlayer: bool\\+Board: *Board\\+Medium, Unbeatable, Minimax
    \nodepart{third}+NewForm(level)\\+Controller() Controller\\+SetLevel(level) error\\+InitializeGame()\\+Click(index) error};
  \node[cls] (board) at (6.4,0) {\centering\textbf{Board}
    \nodepart{second}+Buttons: [9]string\\+Enabled: bool\\+Events: []Event
    \nodepart{third}+SetText(i, mark, who)\\+ClearTexts()\\+ShowMessage(text, caption)\\+TakeEvents() []Event};
  \node[cls] (ctrl) at (0,-4.6) {\centering$\ll$interface$\gg$ \textbf{Controller}
    \nodepart{second}---
    \nodepart{third}+InitializeData()\\+InitializeControls(b)\\+PlayUser(b, i) bool\\+PlayNPC(b)\\+Flag() string\\+Weights() [9]int};
  \node[cls] (strong) at (6.4,-5.2) {\centering\textbf{UnbeatableAIController}
    \nodepart{second}-currentPlayerFlag: Player\\+MoveCount: int\\+CurrentGameState: GameState\\+BoardState: [9]Player\\+ArrayMoves: [9]int\\-heuristicSum: int
    \nodepart{third}+InitializeGameData()\\+PlayUser(b, i) bool\\+PlayNPC(b)\\+CheckWinner(b) GameState\\-heuristicNumericalMethod() int\\-pick(a, b, d) int\\-randomFreeAmong(...) (int, bool)};
  \node[cls] (weak) at (-3.0,-9.8) {\centering\textbf{MediumAIController}\\(Paul Moreno -- implementado)
    \nodepart{second}MoveCount, BoardState, ArrayMoves
    \nodepart{third}InitializeData / CheckWinner / PlayUser / PlayNPC\\HeuristicNumericalMethod};
  \node[pend] (minmax) at (3.0,-9.8) {\centering\textbf{MinimaxAIController}\\(Diana Guerra -- pendiente)
    \nodepart{second}board: [3][3]byte \nodepart{third}GetBestMove / Minimax / EvaluateBoard};
  \draw[->] (api) -- node[above]{usa} (form);
  \draw[{Diamond[open]}-] (form) -- node[above]{1} (board);
  \draw[{Diamond[open]}-] (form) -- node[right]{1..*} (ctrl);
  \draw[dashed,-{Triangle[open]}] (strong) -- (ctrl);
  \draw[dashed,-{Triangle[open]}] (weak.north) -- (ctrl.south);
  \draw[dashed,-{Triangle[open]}] (minmax.north) -- (ctrl.south);
  \draw[->] (strong) -- node[right]{modifica} (board);
\end{tikzpicture}}
\caption{Diagrama de clases del backend}
\end{figure}

% =====================================================================
\section{Metodología SCRUM + XP}
% =====================================================================
Se combinó \textbf{SCRUM} para organizar el trabajo en \emph{sprints} y
\textbf{XP} para las prácticas de ingeniería.

\subsection{Roles}
\begin{table}[H]
\centering
\caption{Roles del equipo}
\begin{tabular}{ll}
\toprule
\textbf{Rol} & \textbf{Responsable}\\
\midrule
Product Owner & Ing. César Villacís Silva, MSc.\\
Scrum Master & Diana Guerra\\
Desarrollo -- UnbeatableAI (Go + Angular) & Germán Cáceres\\
Desarrollo -- MediumAI & Paul Moreno\\
Desarrollo -- UnbeatableAI $\rightarrow$ MinMax & Diana Guerra\\
\bottomrule
\end{tabular}
\end{table}

\subsection{Product Backlog}
\begin{longtable}{p{1.1cm}p{8.6cm}p{1.2cm}p{1.6cm}}
\caption{Historias de usuario}\\
\toprule
\textbf{HU} & \textbf{Historia} & \textbf{Pts} & \textbf{Sprint}\\
\midrule\endfirsthead
\toprule
\textbf{HU} & \textbf{Historia} & \textbf{Pts} & \textbf{Sprint}\\
\midrule\endhead
HU01 & Como usuario quiero ver un tablero 3$\times$3 para poder jugar. & 3 & 1\\
HU02 & Como usuario quiero marcar una casilla y que la JCC responda. & 5 & 1\\
HU03 & Como usuario quiero enfrentar a una IA que gane y bloquee (UnbeatableAI). & 8 & 1\\
HU04 & Como usuario quiero saber quién ganó o si hubo empate. & 3 & 1\\
HU05 & Como usuario quiero elegir el modo de juego. & 3 & 2\\
HU06 & Como usuario quiero animaciones y sonidos para que el juego sea atractivo. & 5 & 2\\
HU07 & Como estudiante quiero ver los pesos y sumas que usa la IA. & 3 & 2\\
HU08 & Como usuario quiero un marcador de partidas. & 2 & 2\\
HU09 & Como usuario quiero jugar contra MediumAI. & 5 & 3\\
HU10 & Como usuario quiero jugar contra UnbeatableAI con MinMax. & 8 & 3\\
\bottomrule
\end{longtable}

\subsection{Sprints}
\begin{description}
  \item[Sprint 1 -- Núcleo del juego.] Port del método numérico a Go
        (\texttt{unbeatable\_ai.go}), tablero (\texttt{board.go}), despacho de clics
        (\texttt{form.go}), API REST y pruebas de los Casos 1--3.
  \item[Sprint 2 -- Experiencia web.] Frontend Angular 22: tablero 3D con
        trazos animados, fondo de partículas, sonidos sintetizados, inspector
        de IA, marcador y pantalla de resultado con confeti.
  \item[Sprint 3 -- MediumAI.] Port del controlador de IA d\'ebil e
        integraci\'on del modo en la interfaz \texttt{Controller}.
\end{description}
\pendiente{Completar y documentar el modo MinMax.}

\subsection{Prácticas XP aplicadas}
\begin{itemize}
  \item \textbf{Diseño simple:} un archivo por modo y una interfaz común.
  \item \textbf{Pruebas primero:} los ejemplos de los PDF se escribieron como
        pruebas (\texttt{unbeatable\_ai\_test.go}) antes de conectar la interfaz.
  \item \textbf{Estándar de código:} \texttt{gofmt}/\texttt{go vet} en Go y Prettier en Angular.
  \item \textbf{Integración continua y entregas pequeñas:} cada incremento se
        sube al repositorio \texttt{github.com/Gpcaceres/videojuegos-grupo3}.
  \item \textbf{Propiedad colectiva:} el \texttt{README} explica cómo
        completar cada modo.
\end{itemize}

\subsection{Diseño}
La arquitectura separa responsabilidades en dos capas:
\begin{itemize}
  \item \textbf{Backend (Go):} \texttt{internal/game} contiene la lógica
        (sin dependencias web) e \texttt{internal/api} expone las rutas.
        \texttt{Board} reemplaza a los controles de Windows Forms
        (\texttt{ListButtons}, \texttt{GroupBox} y \texttt{MessageBox}), de modo
        que el código de los controladores queda casi línea por línea igual al
        original.
  \item \textbf{Frontend (Angular):} componentes \emph{standalone} con
        \emph{signals}, sin \texttt{zone.js}. \texttt{GameStore} reproduce los
        eventos del backend con pausas para animarlos y no decide ninguna jugada.
\end{itemize}

\begin{table}[H]
\centering
\caption{API REST}
\begin{tabular}{lll}
\toprule
\textbf{Método} & \textbf{Ruta} & \textbf{Equivale a (original)}\\
\midrule
GET & /api/modes & Elegir qué aplicación (variante) abrir\\
POST & /api/games & Abrir el formulario (\texttt{frmTicTacToe\_Load})\\
POST & /api/games/\{id\}/moves & \texttt{btnTicTacToeN\_Click}\\
POST & /api/games/\{id\}/new & \texttt{btnNew\_Click}\\
PUT & /api/games/\{id\}/level & Cambiar de variante\\
\bottomrule
\end{tabular}
\end{table}

\subsection{Desarrollo}
La correspondencia entre el original y el port es directa:

\begin{table}[H]
\centering
\caption{Correspondencia C\# $\rightarrow$ Go}
\small
\begin{tabular}{ll}
\toprule
\textbf{C\# (original)} & \textbf{Go}\\
\midrule
\texttt{ListButtons[i].Text} & \texttt{Board.Buttons[i]}\\
\texttt{grbTicTacToe.Enabled} & \texttt{Board.Enabled}\\
\texttt{MessageBox.Show(...)} & \texttt{Board.ShowMessage(...)} $\rightarrow$ evento\\
\texttt{enum Player} / \texttt{enum GameState} & \texttt{type Player int} / \texttt{type GameState int}\\
\texttt{mBoardState} (\texttt{Player[]}) & \texttt{BoardState [9]Player}\\
\texttt{mArrayMoves} (\texttt{int[]}) & \texttt{ArrayMoves [9]int}\\
\texttt{CENTER\_INDEX}, \texttt{CORNER\_/SIDE\_INDICES} & \texttt{CenterIndex}, \texttt{Corner/SideIndices}\\
\texttt{new Random().Next(0, 9)} & \texttt{rand.Intn(9)}\\
\texttt{mHeuristicSum} (miembro) & \texttt{heuristicSum} (se conserva entre turnos)\\
\texttt{static bool ControlPlayer} & \texttt{Form.ControlPlayer} (por sesión)\\
\bottomrule
\end{tabular}
\end{table}

La única diferencia deliberada es la protección de los \texttt{do/while}
aleatorios: si no queda ninguna casilla del grupo pedido, el original entra
en un bucle infinito. En el servidor esto bloquearía la petición, así que se
continúa con la jugada heurística (RNF07).

\subsection{Pruebas}
\begin{table}[H]
\centering
\caption{Resultados de \texttt{go test ./...}}
\small
\begin{tabular}{p{6.2cm}p{6.2cm}c}
\toprule
\textbf{Prueba} & \textbf{Verifica} & \textbf{Resultado}\\
\midrule
TestFirstMoveTakesCenter & La JCC toma el centro en su primera jugada & PASS\\
TestFirstMoveCornerWhenCenterTaken & Si el centro está ocupado, toma una esquina (50 repeticiones) & PASS\\
TestCase1Blocks & Caso 1: bloquea al usuario (suma 2) & PASS\\
TestCase2Wins & Caso 2: completa la línea y gana (suma 6) & PASS\\
TestTakesImmediateWin & Prioriza ganar en la diagonal & PASS\\
TestAdjacentSidesCornerOrder & Orden oficial de las esquinas (1+3$\to$0 antes que 1+5$\to$2) & PASS\\
TestSideRuleNeverPicksThree & Con suma 5 y JCC en el centro, elige lado 1, 5 o 7, nunca 3 (200 repeticiones) & PASS\\
TestOriginalLogicExhaustive & Recorre todas las partidas posibles & PASS\\
\bottomrule
\end{tabular}
\end{table}

Además se verificó la aplicación completa en el navegador (front + back):
la JCC responde en el centro, el resultado se muestra al terminar y los modos
pendientes aparecen bloqueados.

\pendiente{agregar capturas de pantalla de la aplicación en \texttt{Img/} (por ejemplo \texttt{captura-juego.png}).}

% =====================================================================
\section{Análisis de Resultados: Métricas de Software}
% =====================================================================

\subsection{Líneas de código}
Se contaron las líneas efectivas de código (sin comentarios ni líneas en blanco):

\begin{table}[H]
\centering
\caption{Líneas de código por proyecto}
\small
\begin{tabular}{lrrrr}
\toprule
\textbf{Proyecto} & \textbf{Archivos} & \textbf{LOC} & \textbf{Coment.} & \textbf{Blanco}\\
\midrule
Referencia del docente, IA Fuerte (C\#) & 3 & 559 & 87 & 56\\
\textbf{Backend Go (este proyecto)} & 9 & 854 & 156 & 116\\
\textbf{Frontend Angular (TS + HTML + SCSS)} & 24 & 2\,179 & 36 & 292\\
\bottomrule
\end{tabular}
\end{table}

\begin{table}[H]
\centering
\caption{Clase de la IA: original frente al port}
\begin{tabular}{lrr}
\toprule
\textbf{Archivo} & \textbf{LOC} & \textbf{Funciones/métodos}\\
\midrule
UnbeatableAIController.cs (docente) & 417 & 15\\
\textbf{unbeatable\_ai.go} & \textbf{223} & \textbf{17}\\
\bottomrule
\end{tabular}
\end{table}

El port en Go reduce un 47\,\% las líneas de la clase de IA frente al C\#.
Las mismas reglas se expresan con un \texttt{switch} sin repetir bloques, y
una función auxiliar \texttt{pick} reemplaza los 16 bloques
\texttt{if/else if} iguales (uno por línea y por suma). Las constantes y
enumeraciones compartidas (\texttt{Player}, \texttt{GameState},
\texttt{WinningLines}) quedan en \texttt{controller.go} para que MediumAI las
reutilice.

\subsection{Otras métricas}
\begin{itemize}
  \item \textbf{Pruebas automatizadas:} 8 de 8 aprobadas; la exploración
        exhaustiva cubre alrededor de 15\,000 partidas en 0,02 s.
  \item \textbf{Tamaño del frontend:} 216 kB sin comprimir (62 kB
        transferidos) en el \emph{build} de producción.
  \item \textbf{Dependencias del backend:} cero (solo librería estándar de Go).
  \item \textbf{Tiempo de respuesta:} menos de 1 ms por jugada en el servidor local.
\end{itemize}

\subsection{Hallazgos sobre la lógica del docente}\label{sec:hallazgos}
Al portar y probar \texttt{UnbeatableAIController.cs} se encontró lo siguiente.
Como la especificación exige conservar la lógica, todo se mantiene igual y se
documenta aquí.

\begin{enumerate}
  \item \textbf{No es 100\,\% imbatible.} La prueba exhaustiva muestra que la
        JCC \textbf{pierde en cerca del 1,1\,\% de los caminos posibles}. Cuando
        ninguna regla aplica, la lógica cae en \texttt{GenerateRandomNPCMove}
        y a veces deja al usuario un \emph{tenedor} (dos amenazas a la vez).
        Ejemplo: X en 0, O en 4, X en 7, O en 2 (al azar), X en 6. La JCC solo
        puede tapar una de las amenazas (0--3--6 y 6--7--8), y el usuario gana.
  \item \textbf{El lado 3 nunca se elige.} En \texttt{PlayNPC}, la condición del
        lado aleatorio compara \texttt{SIDE\_INDICES[0]} dos veces
        (\texttt{index == SIDE\_INDICES[0] || index == SIDE\_INDICES[0] || ...}),
        así que el lado 3 queda fuera. Se reprodujo tal cual y la prueba
        \texttt{TestSideRuleNeverPicksThree} lo verifica.
  \item \textbf{Posible bloqueo.} Los \texttt{do/while} aleatorios no tienen
        salida si no queda ninguna casilla válida del grupo. En las cerca de
        15\,000 partidas exploradas eso \textbf{no ocurrió nunca}, pero el
        servidor tiene una protección (RNF07) por si acaso.
  \item \textbf{Regla del centro sin efecto.} A partir de la segunda jugada,
        \texttt{PlayNPC} revisa si el centro está libre; nunca lo está a esa
        altura, porque la primera jugada de alguno de los dos ya lo ocupó.
\end{enumerate}

El primer punto es justamente la motivación de la variante
\textbf{UnbeatableAI $\rightarrow$ MinMax}, que explora el árbol completo de
jugadas y no puede caer en un tenedor.

\pendiente{Agregar m\'etricas de MinMax (LOC, pruebas y porcentaje de derrotas) cuando est\'e implementado.}

% =====================================================================
\section{Conclusiones y Recomendaciones}
% =====================================================================
\subsection{Conclusiones}
\begin{itemize}
  \item Se migró \texttt{UnbeatableAIController} (versión IA Fuerte del
        docente) de C\# a Go sin alterar su lógica. Las pruebas reproducen los
        Casos 1, 2 y 3 del material de clase.
  \item Se integr\'o MediumAI con prioridad al centro, reglas de suma 6 y 2,
        y selecci\'on aleatoria como respaldo, conservando la l\'ogica del docente.
  \item Separar la l\'ogica (backend) de la presentaci\'on (frontend) permiti\'o
        sustituir Windows Forms por una interfaz web animada sin tocar el
        algoritmo.
  \item El método numérico basado en pesos 1/3/0 es simple y eficiente, pero
        no es estrictamente imbatible: su caída al azar deja un 1,1\,\% de
        caminos perdedores.
\end{itemize}

\subsection{Recomendaciones}
\begin{itemize}
  \item Implementar la variante MinMax para garantizar una IA realmente imbatible.
  \item Agregar las reglas de ``bloquear tenedor'' si en el futuro se permite
        modificar la heurística.
  \item Desplegar el backend en un servicio en la nube (por ejemplo Cloud Run)
        y el frontend como sitio estático.
\end{itemize}

% =====================================================================
\begin{thebibliography}{9}
\bibitem{ieee830} IEEE, \emph{IEEE Recommended Practice for Software Requirements Specifications}, IEEE Std 830-1998.
\bibitem{villacis} C. Villacís, W. Fuertes, A. Bustamante et al., \emph{Optimización del Juego del Tres en Raya con Niveles de Dificultad, Utilizando Técnicas Heurísticas de Inteligencia Artificial}, Universidad de las Fuerzas Armadas -- ESPE.
\bibitem{go} The Go Authors, \emph{The Go Programming Language Documentation}, \url{https://go.dev/doc/}.
\bibitem{angular} Google, \emph{Angular Documentation}, \url{https://angular.dev}.
\bibitem{scrum} K. Schwaber y J. Sutherland, \emph{La Guía de Scrum}, 2020.
\bibitem{xp} K. Beck, \emph{Extreme Programming Explained: Embrace Change}, Addison-Wesley, 2004.
\end{thebibliography}

\end{document}
